%! Boxplots and Comparing Distributions — Unit 1, Lesson 1.6 — Experience & Formalize
%! (Activity · QuickNotes · Application) (Answer Key)
% EFFL component, Math Medic sizing: 12pt body, OPEN white space after every question.
% Page budget: THREE pages. Activity <=2pp · QuickNotes 1/2pp · Application 1/2-1pp.
% Check Your Understanding is NOT in this component: those items were moved to the end of
% the packet and are now the lesson's SCORED homework (see ../homework/).
%
% PAGE BUDGET NOTE: an earlier draft ran to 4pp with the Application stranded alone on p.4.
% The page was recovered from QUESTION PROSE and STRUCTURAL SPACING, never from write space:
% every \answerspace is at its usable floor (1.4-2.6cm) and must not be shaved further.
% What was cut: merged lead-in paragraphs, one line of throat-clearing per sub-question stem,
% the doubled header row on the 1b table, \arraystretch 1.3 -> 1.2, itemsep 6pt -> 4pt,
% tighter TikZ vertical extents and scales, and connective prose in QuickNotes (every
% fill-in blank was kept). Any of that re-expanding pushes the Application back onto p.4.
%
% THIS LESSON MERGES CED TOPICS 1.8 AND 1.9, so the activity holds only TWO problems. The
% comparison reps (Topic 1.9) are carried by the Application and the homework.
%
% SPOILER RULE: the Activity never says five-number summary, boxplot, box, whisker, modified
% boxplot, or parallel boxplots. The displays are called "a summary picture", "the rectangle",
% and "the lines". Median, Q1, Q3, IQR, range, and the 1.5 x IQR rule WERE formalized in
% Lesson 1.5 and are prior knowledge here; shape and outlier came from Lesson 1.4.
%
% NO SKETCH QUESTIONS (hard constraint): every display is pre-drawn. Students "construct" only
% by filling in the five-number table in 1a, never by drawing an axis or a box.
%
% THE DATA (problem 1: free throws made out of 100, twenty players):
%   30 31 31 32 33 | 36 36 37 38 38 | 41 43 45 47 49 | 52 56 60 66 74
%   min 30 · Q1 34.5 · median 39.5 · Q3 50.5 · max 74
%   Section widths: 4.5, 5.0, 11.0, 23.5 — the right section is over five times the left —
%   yet each holds exactly FIVE of the twenty. That count is the discovery (EK UNC-1.L.2).
%   Every boundary falls inside a gap in the data (33|36, 38|41, 49|52), so the counting in
%   1b is unambiguous: 5, 5, 5, 5. PROTECTED — do not soften the width contrast.
%
% THE DATA (problem 2, travel minutes): Route A 18/21/23/26/31 (IQR 5, no outliers);
%   Route B 16/20/25/32/58 (IQR 12; upper cut-off 32 + 1.5(12) = 50, so 58 is set apart and
%   the line stops at 40, the most extreme value that is NOT unusually far).
% GENERATED FROM experience/main.tex — the blank with answers filled in, so every
% \answerspace height, every \boxguard, every scale=, every \vspace, every itemsep and the
% \bxp macro are identical to it. NO teachernote here — that lives in the lesson plan.
\documentclass[12pt]{article}
\usepackage{apstats-article}
\usepackage{apstats-key}   % pulls in -boxes; do NOT also load -boxes
\newcommand{\answerspace}[2]{\par\nopagebreak\noindent\begin{minipage}[t][#1][t]{\linewidth}%
  \color{keyred}\bfseries #2\end{minipage}\par}
% \bxp{y-center}{half-height}{min}{Q1}{median}{Q3}{max} — every display in this lesson is
% PRE-DRAWN with it. Authored byte-identically in the blank and the key.
\newcommand{\bxp}[7]{%
  \draw[thick, navy] (#3,#1) -- (#4,#1);
  \draw[thick, navy] (#6,#1) -- (#7,#1);
  \draw[thick, navy] (#3,{#1-#2}) -- (#3,{#1+#2});
  \draw[thick, navy] (#7,{#1-#2}) -- (#7,{#1+#2});
  \draw[thick, navy, fill=skymid] (#4,{#1-#2}) rectangle (#6,{#1+#2});
  \draw[very thick, navy] (#5,{#1-#2}) -- (#5,{#1+#2});}

\begin{document}

\pageheader{Unit 1, Lesson 1.6}{Experience \& Formalize: Twenty Free Throws --- Answer Key}
% NAMESTRIP: no \namedateperiod here — mirrors the blank.

\vspace{0.02in}

% ── Part 1: the Activity ─────────────────────────────────────────────────────
\begin{headlinebox}{sky}
\textbf{Twenty Free Throws.} A coach recorded how many free throws each of the twenty players
on the squad made out of 100 attempts this week. Work both problems with your group using only
what you already know --- and in problem~1, \textbf{count} before you conclude.
\end{headlinebox}

\vspace{0.04in}

\begin{scenariobox}[1. The Same Twenty, Twice]{navy}
Every player is one dot above the axis. Below the axis is a summary picture somebody drew from
those same twenty numbers.

\vspace{3pt}
\begin{center}
\begin{tikzpicture}[scale=0.8]
  \foreach \x in {0,1,2,3,4,5,6,7,8,9} {
    \draw[linegray, very thin] (\x,0.02) -- (\x,0.55);
    \draw[linegray, very thin] (\x,-1.70) -- (\x,-0.85);}
  \foreach \x/\n in {0/1, 0.2/2, 0.4/1, 0.6/1, 1.2/2, 1.4/1, 1.6/2, 2.2/1, 2.6/1, 3.0/1,
                     3.4/1, 3.8/1, 4.4/1, 5.2/1, 6.0/1, 7.2/1, 8.8/1} {
    \foreach \k in {1,...,\n} { \fill[navy] (\x,{0.2*\k}) circle (0.075); }}
  \draw[->, thick] (-0.35,0) -- (9.4,0);
  \foreach \x/\lab in {0/30, 1/35, 2/40, 3/45, 4/50, 5/55, 6/60, 7/65, 8/70, 9/75} {
    \draw (\x,-0.08) -- (\x,0.08); \node[below, font=\scriptsize] at (\x,-0.12) {\lab};}
  \node[font=\scriptsize] at (4.5,-0.55) {Free throws made (out of 100)};
  \bxp{-1.25}{0.36}{0}{0.9}{1.9}{4.1}{8.8}
\end{tikzpicture}
\end{center}
\vspace{2pt}

\normalsize
\begin{enumerate}[label=\textbf{\alph*.}, leftmargin=1.8em, itemsep=4pt]
  \item Read the twenty values off the top picture and fill in these five numbers.

        \vspace{2pt}
        \renewcommand{\arraystretch}{1.2}
        \begin{tabularx}{\linewidth}{@{} Y Y Y Y Y @{}}
        \rowcolor{sky}
        \textbf{Least} & \textbf{$Q_1$} & \textbf{Median} & \textbf{$Q_3$} & \textbf{Greatest} \\
        \hline
        \ans{30} & \ans{34.5} & \ans{39.5} & \ans{50.5} & \ans{74} \\
        \end{tabularx}

        \vspace{4pt}
        Where does each of those five show up in the summary picture?
        \answerspace{1.4cm}{30 and 74 are the far ends of the two lines, 34.5 and 50.5 are the
        edges of the rectangle, and 39.5 is the line inside it.}
  \item Those five cut the picture into four sections. \textbf{Count} the dots in each.

        \vspace{2pt}
        \renewcommand{\arraystretch}{1.2}
        \begin{tabularx}{\linewidth}{@{} l Y Y Y Y @{}}
        \rowcolor{sky}
        \textbf{Range} & 30 to 34.5 & 34.5 to 39.5 & 39.5 to 50.5 & 50.5 to 74 \\
        \hline
        \textbf{Players} & \ans{5} & \ans{5} & \ans{5} & \ans{5} \\
        \end{tabularx}

        \vspace{4pt}
        The right section is over five times as wide as the left. In one sentence, what do your
        counts say about it?
        \answerspace{1.6cm}{The same number --- five each. The wide section's five are spread
        from 52 to 74; the narrow section's five are packed into 30 to 33.}
  \item A classmate says, \emph{``the wide part is where most of the shots landed.''} Use your
        counts to say what is wrong --- and what the width really tells you.
        \answerspace{1.6cm}{It holds five players, the same as every other section. Its width
        says those five totals are far apart, not that most shots landed there.}
  \item Name one thing the top picture shows that the summary picture cannot.
        \answerspace{1.4cm}{The individual values --- it cannot give back that two players made
        31, or show the tight clump from 30 to 38.}
\end{enumerate}
\end{scenariobox}

\boxguard[16]
\begin{scenariobox}[2. Two Bus Routes]{navy}
A commuter timed the same morning trip on two bus routes. The five numbers for each are given
below the pictures.

\vspace{3pt}
\begin{center}
\begin{tikzpicture}[scale=0.85]
  \foreach \x in {0,0.8,1.6,2.4,3.2,4.0,4.8,5.6,6.4,7.2} {
    \draw[linegray, very thin] (\x,-0.10) -- (\x,1.80);}
  \bxp{1.40}{0.30}{0.48}{0.96}{1.28}{1.76}{2.56}
  \bxp{0.50}{0.30}{0.16}{0.80}{1.60}{2.72}{4.00}
  \fill[navy] (6.88,0.50) circle (0.085);
  \node[font=\bfseries\scriptsize, anchor=east] at (-0.22,1.40) {Route A};
  \node[font=\bfseries\scriptsize, anchor=east] at (-0.22,0.50) {Route B};
  \draw[->, thick] (-0.20,-0.10) -- (7.75,-0.10);
  \foreach \x/\lab in {0/15, 0.8/20, 1.6/25, 2.4/30, 3.2/35, 4.0/40, 4.8/45, 5.6/50, 6.4/55,
                       7.2/60} {
    \draw (\x,-0.18) -- (\x,-0.02); \node[below, font=\scriptsize] at (\x,-0.22) {\lab};}
  \node[font=\scriptsize] at (3.6,-0.72) {Travel time (minutes)};
\end{tikzpicture}
\end{center}

\renewcommand{\arraystretch}{1.2}
\begin{center}
\begin{tabularx}{0.92\linewidth}{@{} l Y Y Y Y Y @{}}
\rowcolor{sky}
\textbf{Route} & \textbf{Least} & \textbf{$Q_1$} & \textbf{Median} & \textbf{$Q_3$} &
\textbf{Greatest} \\
\hline
\textbf{A} & 18 & 21 & 23 & 26 & 31 \\
\textbf{B} & 16 & 20 & 25 & 32 & 58 \\
\end{tabularx}
\end{center}
\vspace{3pt}

\normalsize
\begin{enumerate}[label=\textbf{\alph*.}, leftmargin=1.8em, itemsep=4pt]
  \item How wide is the middle half of each route's times? \quad Route A: \ans{$26-21=5$ min} \quad
        Route B: \ans{$32-20=12$ min}
  \item Which route is typically faster? Which is less predictable? How can you tell?
        \answerspace{1.6cm}{Route A is faster --- median 23 minutes against 25. Route B is far
        less predictable: its middle half is 12 minutes wide against Route A's 5.}
  \item One Route B rider is drawn as a lone dot. Using the middle-half width from (a), work out
        how far past the rectangle he sits --- then say why whoever drew this set him apart
        instead of stretching the line out to reach him.
        \answerspace{2.0cm}{Route B's middle half is 12 minutes wide, and he sits $58-32=26$
        minutes past the right edge --- over two of those widths. The cut-off is
        $32 + 1.5(12) = 50$, and $58 > 50$, so he is set apart.}
  \item What would you tell a commuter who must not be late? Defend it from \emph{both} pictures.
        \answerspace{1.4cm}{Route A. Its slowest trip here, 31 min, beats Route B's $Q_3$ of 32
        --- a quarter of B's mornings are worse than all of A's.}
\end{enumerate}
\end{scenariobox}

% ── Part 2: QuickNotes (the debrief fills this) — target: half a page ────────
% Every fill-in blank from the original draft was kept (18 of them); only the connective
% prose around them was compressed. Do not re-expand the sentences.
\boxguard[14]
\begin{tcolorbox}[enhanced, colback=sky, colframe=navy, boxrule=0.8pt, arc=2mm,
    left=3mm, right=3mm, top=1.5mm, bottom=1.5mm, breakable,
    fonttitle=\bfseries\small\color{white}, title={QuickNotes: Boxplots and Comparing Distributions},
    attach boxed title to top left={yshift=-2mm, xshift=4mm},
    boxed title style={colback=navy, colframe=navy, arc=1.5mm}]
\small
\begin{itemize}[leftmargin=1.1em, itemsep=3pt]
  \item \textbf{Five-number summary:} minimum, \ans{$Q_1$}, median, \ans{$Q_3$}, maximum.
  \item \textbf{Boxplot:} those five drawn to scale. \textbf{Box} $= Q_1$ to $Q_3$, the middle
        \ans{50}\%; the line inside is the \ans{median}. \textbf{Whiskers} reach the most
        extreme value that is \ans{not} an outlier; an outlier gets its own \ans{symbol}
        (a \textbf{modified boxplot}).
  \item \textbf{THE RULE:} each section holds \ans{25}\% of the data. Width shows
        \ans{spread} --- never \ans{count}.
  \item \textbf{Parallel (side-by-side) boxplots:} one plot per group on a shared \ans{axis}.
  \item \textbf{Comparing:} \ans{center}, \ans{variability}, unusual features --- with a comparative
        word (\emph{higher than}) and always \ans{in context}.
  \item \textbf{Mean vs.\ median:} symmetric, close; skewed right, mean to the \ans{right};
        skewed left, to the \ans{left}.
  \item \textbf{A boxplot cannot show} \ans{gaps}, \ans{clusters}, or a second \ans{peak}.
\end{itemize}
\end{tcolorbox}

% ── Part 3: the Application (worked together, ~6 min) ────────────────────────
\boxguard[14]
\begin{notesbox}{Application: Two Science Sections}
We do this one together. Two sections of the same course, one test.

\vspace{2pt}
\begin{center}
\begin{tikzpicture}[scale=0.85]
  \foreach \x in {0,1.4,2.8,4.2,5.6,7.0} {
    \draw[linegray, very thin] (\x,-0.15) -- (\x,1.75);}
  \bxp{1.35}{0.30}{1.12}{3.36}{4.20}{5.04}{6.16}
  \bxp{0.50}{0.30}{1.40}{2.52}{3.08}{4.76}{6.72}
  \node[font=\bfseries\scriptsize, anchor=east] at (-0.22,1.35) {Period 2};
  \node[font=\bfseries\scriptsize, anchor=east] at (-0.22,0.50) {Period 5};
  \draw[->, thick] (-0.20,-0.15) -- (7.55,-0.15);
  \foreach \x/\lab in {0/50, 1.4/60, 2.8/70, 4.2/80, 5.6/90, 7.0/100} {
    \draw (\x,-0.23) -- (\x,-0.07); \node[below, font=\scriptsize] at (\x,-0.27) {\lab};}
  \node[font=\scriptsize] at (3.5,-0.78) {Test score (points)};
\end{tikzpicture}
\end{center}

\renewcommand{\arraystretch}{1.2}
\begin{center}
\begin{tabularx}{0.92\linewidth}{@{} l Y Y Y Y Y @{}}
\rowcolor{sky}
\textbf{Section} & \textbf{Min} & \textbf{$Q_1$} & \textbf{Median} & \textbf{$Q_3$} &
\textbf{Max} \\
\hline
\textbf{Period 2} & 58 & 74 & 80 & 86 & 94 \\
\textbf{Period 5} & 60 & 68 & 72 & 84 & 98 \\
\end{tabularx}
\end{center}
\vspace{3pt}

\begin{enumerate}[label=\textbf{\alph*.}, leftmargin=1.8em, itemsep=4pt]
  \item Median: Period 2 \ans{80} \quad Period 5 \ans{72} \qquad
        IQR: Period 2 \ans{12} \quad Period 5 \ans{16}
  \item Which section is skewed, and which way? How does the picture tell you?
        \answerspace{1.6cm}{Period 2 is skewed left --- its lower whisker is long and its upper
        half short. Period 5 is skewed right, with its upper half stretched.}
  \item \textbf{One paragraph}, in context: center, variability, unusual features --- with
        comparative language.
        \answerspace{2.6cm}{Period 2's scores were typically higher than Period 5's --- median 80
        points against 72. Period 5's were more variable: a middle half 16 points wide against
        12, and a wider span. Period 2 trails toward low scores, Period 5 toward high; neither
        shows an outlier.}
\end{enumerate}
\end{notesbox}

\end{document}
